%!TEX program = xelatex
% 完整编译: xelatex -> biber/bibtex -> xelatex -> xelatex
\documentclass[lang=cn,a4paper,newtx,code=minted]{elegantpaper}

\title{基于多智能体协同与可视化的中国古典名著知识图谱构建}
\author{杨晨 \and 杨智超 \and 黄金辉}
\institute{大数据学院 \quad 应用统计}
\date{\zhdate{2025/12/28}}

\usepackage{array}
\newcommand{\ccr}[1]{\makecell{{\color{#1}\rule{1cm}{1cm}}}}
\addbibresource{../src/reference.bib} % 参考文献，不要删除


\begin{document}

\maketitle

\begin{abstract}
中国古典名著中交织着复杂的人物关系与文化线索，传统阅读方式难以从全局把握其社会结构。本文以《红楼梦》《三国演义》《水浒传》《西游记》为对象，构建统一 Schema 下的知识图谱，并提出人参与（Human-in-the-loop）的抽取流程：由大语言模型生成候选三元组，在 Schema 与证据约束下进行二次验真，随后结合别名映射与关系归一化完成数据治理。在该流程下，我们得到包含 3,504 个实体与 6,340 条关系的知识图谱，并将原始抽取阶段 1400+ 种长尾关系收敛到 100 种以内。基于 Neo4j，我们实现了图谱可视化与路径检索式问答的原型系统，并进一步探索图谱检索增强生成与图算法工具化调用等应用形态。

\keywords{知识图谱, 大语言模型, 多智能体协同, 古典名著}
\end{abstract}

\section{引言}

中国古典名著不仅提供了丰富的叙事文本，也天然蕴含着复杂的社会网络与文化线索：人物之间的亲属、师承、结义、阵营与冲突共同构成跨章节的关系网。将这类关系结构化为知识图谱（Knowledge Graph）能够把“人物—关系—证据”以可计算的形式组织起来，一方面帮助读者从“故事”上升到“结构”理解作品，另一方面为数字人文中的关系分析、情节检索与知识库问答提供基础设施。

与现代新闻或百科文本相比，古典小说的语言表达更含蓄，称谓与指代更灵活，同一人物往往存在多种别名与尊称；同时，叙事中大量动作性描写会带来关系类型的长尾膨胀。若直接依赖大模型端到端抽取，常见风险包括：将语境暗示误当作事实关系，或在长文本场景下生成缺乏出处的内容。上述因素共同提高了“抽得到、抽得准、抽得统一”的难度。

为此，我们以《红楼梦》《三国演义》《水浒传》《西游记》四部名著为对象，构建统一 Schema 下的知识图谱，并设计人参与（Human-in-the-loop）的抽取流程：模型负责生成候选三元组，规则约束与二次验真负责过滤与校正，最终形成可追溯、可复用的数据资产。基于该图谱，我们实现了面向展示与检索的原型系统，支持图谱浏览、路径检索与可解释问答等交互能力。

\section{数据集与知识图谱的构造}

数据源选取四大名著的通行版本。为适配自动化处理，我们对文本做了两步预处理：一是清洗排版噪声并剔除非正文内容（序言、后记等）；二是按章节标题切分为稳定的叙事单元，并在章节内部按段落边界进一步分块，以尽量保留对话与指代信息。所有文本块均保留章节来源，用于后续三元组的证据回溯与展示。

在 Schema 设计上，我们采用“核心统一、可扩展”的思路。实体侧以人物为中心，并补充组织、地点、物品、事件、文本等类型，用于覆盖名著叙事中的常见对象；关系侧将大量细碎动词归并为更少的核心关系集合，并保留必要的语义区分，以兼顾阅读可解释性与后续图计算的稳定性。为保证跨书目一致性，我们维护了两类映射：
\begin{itemize}
    \item \textbf{实体别名映射：} 将同一人物的不同称谓（如“宝玉”“宝二爷”）对齐到同一规范名，避免实体碎片化。
    \item \textbf{关系归一化映射：} 将低频动作型关系收敛到更少的核心关系类型，降低长尾稀疏性并提升统计稳定性。
\end{itemize}

综合以上设计，图谱构造流程可以概括为：\textbf{文本切分}（保留来源）$\rightarrow$ \textbf{候选抽取}（输出三元组与证据）$\rightarrow$ \textbf{二次验真}（剔除无证据/指代不清）$\rightarrow$ \textbf{清洗归一}（实体对齐与关系收敛）$\rightarrow$ \textbf{入库与服务}（写入 Neo4j，供可视化与检索调用）。

完成抽取与清洗后，我们得到的图谱包含 3,504 个实体与 6,340 条关系，规模分布如表 \ref{tab:kg_stats} 所示。

\begin{table}[htbp]
    \centering
    \caption{知识图谱构建数据统计}
    \label{tab:kg_stats}
    \begin{tabular}{lccc}
        \hline
        \textbf{书名} & \textbf{实体总数} & \textbf{关系总数} & \textbf{原始关系数} \\
        \hline
        \textbf{红楼梦} & 849 & 1,720 & 200+ \\
        \textbf{三国演义} & 1,097 & 2,055 & 500+ \\
        \textbf{水浒传} & 1,055 & 1,803 & 400+ \\
        \textbf{西游记} & 536 & 762 & 351 \\
        \hline
        \textbf{合计} & \textbf{3,504} & \textbf{6,340} & \textbf{1400+} \\
        \hline
    \end{tabular}
\end{table}

\section{任务定义}

本项目围绕“知识抽取”与“下游应用”两条主线展开。给定原始文本块 $x$（包含书目与章节来源），知识抽取任务需要输出结构化结果 $G=(V,E)$：其中 $V$ 为实体集合，$E$ 为关系集合。我们以三元组 $(h,r,t)$ 表示一条关系，并为每条三元组附加两类关键属性：
\begin{itemize}
    \item \textbf{类型约束：} 为实体标注类型（人物/组织/地点/物品/事件/文本等），并要求关系满足 Schema 约束（例如“居住”指向地点）。
    \item \textbf{证据约束：} 为每条关系提供可在原文中定位的证据片段（原句或连续子串），用于验真与解释展示。
\end{itemize}

在下游应用上，我们面向 demo 设计了两类能力：其一是\textbf{图谱浏览与检索}（按实体类型与局部子图查看、邻域探索）；其二是\textbf{可解释问答}：将自然语言问题映射为图数据库查询（Cypher），在返回答案的同时给出推理路径与文本证据。对于需要多跳推理的问题，我们进一步探索\textbf{图谱增强生成}：先从图谱检索相关子图并序列化为背景知识，再进行生成式回答，以增强事实约束。

\section{方法}

方法部分重点回答两个问题：\textbf{如何从文本稳定地产生“有证据”的三元组}，以及\textbf{如何在跨书目场景下保持口径一致}。为降低幻觉与误判风险，我们采用“生成—验真—归一”的流水线，并在关键节点引入人工抽检与规则约束。

\textbf{（1）候选抽取：Generator。} 在抽取阶段，我们为模型注入 Schema（实体类型、关系白名单与类型一致性要求），并要求输出同时包含三元组与证据片段。为了提升复杂关系识别能力，提示中允许模型先进行简短推理再作答（Chain-of-Thought 思路），但最终只保留可由原文证据支持的结果。

\textbf{（2）二次验真：Critic。} 生成结果容易包含“合理但无据”的关系。我们设置评审模型对每条候选关系进行核验：若证据片段无法在原文中定位，或指代对象不清晰，则直接剔除；对边界含糊的实体提及，优先要求回到文本补充上下文后再判定。该阶段的目标不是最大化召回，而是确保最终图谱具备可追溯性。

\textbf{（3）清洗与归一：} 跨书目构建时，主要噪声来自“同人异名”和“关系长尾”。我们维护别名映射表完成实体对齐，并以关系映射表将动作型长尾关系收敛到较少的核心关系集合，从而提升图谱在统计、检索与图算法上的稳定性。对难以自动归一的关系类型，采取人工小规模审校后再扩展映射。

\section{实验}

实验部分聚焦知识抽取任务，重点评估两点：\textbf{（i）抽取结果是否可被证据支撑}，以及\textbf{（ii）归一化是否有效抑制长尾稀疏}。由于古典文本缺少现成标注集，我们采用抽样人工核验与对比实验相结合的方式：在四部书中按章节分层抽样，分别统计“仅 Generator”与“Generator+Critic+归一化”两种设置下的结果差异。为保证结论可复核，抽样时记录文本来源、证据片段、三元组与判定理由。

在评测指标上，我们记录：\textbf{关系证据通过率}（关系是否能在文本中定位到证据）、\textbf{实体对齐一致性}（同一实体是否被统一命名）、以及\textbf{关系类型数}（归一化前后关系长尾规模）。其中，关系类型规模的变化可以直接由构建结果统计给出：原始关系类型 1400+，归一化后收敛到 100 以内（见表 \ref{tab:kg_stats}）。

\begin{table}[htbp]
    \centering
    \caption{知识抽取任务评测记录}
    \label{tab:ie_eval}
    \begin{tabular}{lccc}
        \hline
        \textbf{设置} & \textbf{关系证据通过率} & \textbf{实体对齐一致性} & \textbf{关系类型数} \\
        \hline
        仅 Generator & 65\% & 80\% & 1400+ \\
        + Critic + 归一化 & 92\% & 96\% & < 100 \\
        \hline
    \end{tabular}
\end{table}

误差分析显示，主要难点来自含蓄表达（借喻与转述）以及跨段落指代（如“其”“他”“那人”）导致的实体边界不稳定；对应的改进方向包括更精细的指代消解与跨块上下文补全机制。

\section{原型系统}

为便于展示与交互，我们将图谱存储在 Neo4j 中，并实现了前后端分离的原型系统：
\begin{itemize}
    \item \textbf{后端}：基于 Python Flask，负责数据接口与查询服务，并将用户请求映射为 Cypher 查询。
    \item \textbf{数据库}：采用 Neo4j 存储图谱数据，支持多跳关系检索与路径查询。
    \item \textbf{前端}：使用 Bootstrap 组织页面布局，并通过 ECharts 绘制关系网络，支持缩放、拖拽与邻域高亮等交互。
\end{itemize}

系统提供面向浏览与检索的基本能力：用户可按实体类型筛选并查看局部子图；点击节点可在侧边栏展示人物属性与关联信息；在问答场景中，系统支持路径检索式查询，除返回答案外，还会在图上高亮推理路径，用于解释“答案从何而来”。

为了便于 demo 展示，我们将典型使用流程整理如下：
\begin{enumerate}
    \item \textbf{图谱浏览：} 在前端选择书目或人物实体，加载其邻域子图；通过缩放、拖拽与邻域高亮快速定位关键角色与关系。
    \item \textbf{节点详情：} 点击节点查看属性信息与相关关系，结合章节来源回溯该关系的文本证据。
    \item \textbf{路径检索问答：} 输入“某人物与某人物是什么关系/如何关联”等问题，后端解析实体与关系意图并生成 Cypher 查询，返回答案与可视化路径。
    \item \textbf{证据与溯源：} 对关键关系展示证据片段与章节出处，支持对结果进行复核与纠错，形成可解释闭环。
\end{enumerate}

在应用探索方面，我们补充了两类与大模型结合的用法：其一是在复杂多跳问题中，先从图谱检索相关子图并序列化为背景知识（Graph-RAG 思路），再进行生成式回答，以增强事实约束；其二是将最短路径、社区发现等图算法封装为可调用工具，在需要时由智能体触发执行，将结构化计算纳入推理流程。

\section{结论}

本文围绕四大名著构建了统一 Schema 的知识图谱，并在抽取流程中引入“生成—验真—清洗”的组合策略，以提升关系的可溯源性与跨书目一致性；同时实现了基于 Neo4j 的原型系统，用于验证图谱在可视化浏览与可解释检索中的使用价值。

后续工作将重点放在三方面：其一，引入更明确的时序信息，将关系与章节/时间线绑定，以支持情节演化分析；其二，完善实体别名与指代消解，提升跨段落、跨章节的实体一致性；其三，扩展到更多古籍文本，并探索中英双语对齐与多模态（插图/影视素材）关联，提升知识图谱的可用范围与传播效果。

\section{人员与工作安排}
本项目采用分工协作的方式推进：以知识抽取与数据治理为主线，配合原型系统实现与展示，最终形成可复现实验与可演示 demo 的完整闭环。三位成员在各自负责模块之外，也参与了联调、结果核验与文档打磨，确保交付物一致且可用。
\begin{itemize}
    \item \textbf{杨晨（组长）}: 负责整体框架设计与工作流搭建，包括“生成—验真—归一”的流程规划、关键模块接口约定与阶段性里程碑推进；同时统筹汇报材料，组织组内讨论与结果复核，保证技术路线与报告口径一致。
    \item \textbf{黄金辉}: 负责前端可视化展示与页面设计，实现图谱浏览、节点详情与交互功能，并配合后端接口完成 demo 联调；同时承担文档整理与排版优化工作，完善系统使用说明与展示流程。
    \item \textbf{杨智超}: 负责知识抽取与清洗脚本开发，包括文本预处理、候选三元组抽取、别名映射与关系归一化等数据治理环节；同时参与抽样核验与错误分析，协助完善映射规则与实验记录。
\end{itemize}

\section{组内互评}
\noindent
本次项目组内互评以“协作态度、任务完成度、交付质量、沟通响应”四个维度为参考。经组内讨论，三位成员对彼此评价一致，整体评价为“很好”，汇总如表 \ref{tab:peer_review} 所示。

\begin{table}[htbp]
    \centering
    \caption{组内互评汇总}
    \label{tab:peer_review}
    \begin{tabular}{p{0.18\textwidth} p{0.18\textwidth} p{0.56\textwidth}}
        \hline
        \textbf{评价人} & \textbf{被评价人} & \textbf{评价（均为“很好”）} \\
        \hline
        杨晨 & 黄金辉 & 前端可视化实现与页面设计完成度高，demo 展示思路清晰，联调响应及时，并能主动补齐文档与展示材料。 \\
        杨晨 & 杨智超 & 抽取与清洗脚本实现扎实，数据处理链路清晰，能快速定位问题并迭代规则，支撑了图谱构建与实验分析。 \\
        黄金辉 & 杨晨 & 整体框架设计与推进有条理，任务拆解清晰，统筹沟通顺畅，汇报组织到位，保证了项目按期完成。 \\
        黄金辉 & 杨智超 & 数据抽取与归一化工作细致，输出稳定可靠，配合接口与数据格式调整迅速，为前端展示提供了高质量数据。 \\
        杨智超 & 杨晨 & 技术路线把控准确，工作流设计清晰，能及时协调资源与解决阻塞问题，对最终报告质量提升帮助明显。 \\
        杨智超 & 黄金辉 & 可视化与交互实现效果直观，能够将图谱能力以 demo 形式清楚呈现，沟通配合积极，推动了系统闭环落地。 \\
        \hline
    \end{tabular}
\end{table}

% \printbibliography[heading=bibintoc, title=参考文献]

\end{document}
